\documentclass[10pt,a4paper]{article}
\usepackage[top=1.5cm, bottom=1.5cm, left=1.6cm, right=1.6cm, headheight=14pt, footskip=18pt]{geometry}
\usepackage[utf8]{inputenc}
\usepackage{amsmath,amsfonts,amssymb}
\usepackage{graphicx}
\usepackage{tikz}
\usepackage{circuitikz}
\usetikzlibrary{calc,patterns}
\usepackage[most]{tcolorbox}
\usepackage{enumitem}
\usepackage{fancyhdr}
\usepackage{booktabs}
\usepackage{tabularx}
\usepackage{colortbl}
\usepackage{hyperref}

% --- Palet Warna Resmi UNIB ---
\definecolor{unibblue}{RGB}{0, 32, 96}
\definecolor{unibgold}{RGB}{197, 160, 89}
\definecolor{unibgreen}{RGB}{34, 112, 60}
\definecolor{boxbg}{RGB}{248, 250, 253}
\definecolor{linegray}{RGB}{210, 215, 225}

% --- Pengaturan Header & Footer ---
\pagestyle{fancy}
\fancyhf{}
\lhead{\footnotesize\textbf{\color{unibblue}Universitas Bengkulu} \textbar\ Program Studi S1 Teknik Elektro}
\rhead{\footnotesize\color{gray}Persamaan Diferensial (Genap 2026)}
\lfoot{\footnotesize\color{gray}Lembar Kerja Mahasiswa (LKM) Berbasis OBE}
\cfoot{\footnotesize\color{gray}Hal. \thepage\ dari 2}
\rfoot{\footnotesize\color{unibblue}\textbf{Minggu 1: Klasifikasi PDB/PDP, Nilai Awal ...}}
\renewcommand{\headrulewidth}{0.6pt}
\renewcommand{\footrulewidth}{0.4pt}

% --- Custom Box untuk Lembar Jawab ---
\newtcolorbox{answerbox}[1]{%
  colback=white,
  colframe=unibblue!70!black,
  coltitle=white,
  fonttitle=\bfseries\scriptsize,
  colbacktitle=unibblue,
  title={#1},
  arc=2pt,
  left=5pt,right=5pt,top=3pt,bottom=3pt,
  boxrule=0.6pt
}

\begin{document}

% ==================== HALAMAN 1 ====================
\noindent\begin{minipage}[c]{0.68\textwidth}%
  {\large\textbf{\color{unibblue}LEMBAR KERJA MAHASISWA (LKM)}}\\[1mm]
  {\normalsize\textbf{Mata Kuliah: Persamaan Diferensial (2 SKS)}}\\[0.5mm]
  {\small\textbf{Topik Minggu 1:} Klasifikasi PDB/PDP, Nilai Awal (IVP), \& Rangkaian RL}%
\end{minipage}%
\hfill%
\begin{minipage}[c]{0.30\textwidth}%
  \begin{tcolorbox}[colback=unibblue!5,colframe=unibblue,boxrule=0.6pt,arc=2pt,left=4pt,right=4pt,top=2pt,bottom=2pt]
    \scriptsize
    \textbf{Nama:} \dotfill\\[1.2mm]
    \textbf{NPM:} \dotfill\\[1.2mm]
    \textbf{Kelas/Klp:} \dotfill\\[1.2mm]
    \textbf{Nilai Final:} \hfill \textbf{/ 100}
  \end{tcolorbox}%
\end{minipage}\par

\vspace{1.5mm}

% Kotak Sub-CPMK & Pedoman
\begin{tcolorbox}[colback=boxbg,colframe=unibgold!90!black,boxrule=0.6pt,arc=2pt,left=5pt,right=5pt,top=3pt,bottom=3pt,title=\textbf{\scriptsize Capaian Pembelajaran (Sub-CPMK 1) \& Pedoman Evaluasi OBE}]
  \scriptsize
  \textbf{Sub-CPMK 1:} Mahasiswa mampu mengidentifikasi, mengklasifikasi, dan memodelkan fenomena dinamika sistem elektrik terpusat ke dalam Persamaan Diferensial Biasa (PDB) orde 1 serta memvalidasinya secara numerik.\\
  \textbf{Alokasi Waktu:} 50--60 Menit \quad\textbar\quad \textbf{Model:} Kolaboratif / Mandiri Terstruktur \quad\textbar\quad \textbf{Bahasa Komputasi:} Julia 1.12+
\end{tcolorbox}

\vspace{1mm}

% Soal C1
\noindent\textbf{1. (C1 - Mengingat \textbar\ Bobot: 10 Poin)}\par\vspace{0.5mm}
{\footnotesize Sebutkan definisi formal Persamaan Diferensial Biasa (PDB) vs Persamaan Diferensial Parsial (PDP), serta jelaskan kriteria linearitas suatu PDB terhadap variabel terikat beserta seluruh turunannya.}
\begin{answerbox}{Ruang Pengerjaan C1 (Konsep Formal \& Klasifikasi Linearitas)}
  \vspace{2.0cm}
\end{answerbox}

\vspace{1mm}

% Soal C2
\noindent\textbf{2. (C2 - Memahami \textbar\ Bobot: 15 Poin)}\par\vspace{0.5mm}
{\footnotesize Tinjau solusi umum $y(t) = C e^{-t/\tau} + y_{\text{ss}}$. Jelaskan mengapa konstanta $C$ hanya dapat ditentukan jika diberikan Masalah Nilai Awal (IVP), serta jelaskan makna fisis konstanta waktu $\tau = L/R$ pada respons transien arus induktor.}
\begin{answerbox}{Ruang Pengerjaan C2 (Intuisi Fisika Transien \& Nilai Awal)}
  \vspace{2.1cm}
\end{answerbox}

\vspace{1mm}

% Soal C3
\noindent\textbf{3. (C3 - Menerapkan \textbar\ Bobot: 20 Poin)}\par\vspace{0.5mm}
\noindent\begin{minipage}[t]{0.60\textwidth}%
  {\footnotesize Pada rangkaian RL di samping, sakelar ditutup pada $t=0$ dengan kondisi awal $i(0) = 0$. Turunkan persamaan diferensial loop arus dari KVK dan tentukan fungsi analitis arus $i(t)$ untuk $t \ge 0$ jika $V_s = 24\text{ V}$, $R = 8\,\Omega$, dan $L = 0{,}4\text{ H}$.}%
\end{minipage}%
\hfill%
\begin{minipage}[t]{0.37\textwidth}%
  \centering
  \resizebox{0.95\linewidth}{!}{%
  \begin{circuitikz}[american, scale=0.65, transform shape]
    \draw (0,0) to[battery2, l=$V_s$, invert] (0,1.8)
          to[switch, l={$t{=}0$}] (1.6,1.8)
          to[R, l=$R$] (3.2,1.8)
          to[L, l=$L$, v=$v_L(t)$] (3.2,0)
          to[short] (0,0);
    \draw[->, >=stealth, thick, unibblue] (1.4,0.8) arc (160:-160:0.35);
    \node at (1.8,0.8) [unibblue, font=\tiny\bfseries] {$i(t)$};
  \end{circuitikz}%
  }%
\end{minipage}\par\vspace{1mm}
\begin{answerbox}{Ruang Pengerjaan C3 (Penurunan Solusi Eksak \& Perhitungan Nilai)}
  \vspace{2.8cm}
\end{answerbox}

\newpage
% ==================== HALAMAN 2 ====================

% Soal C4
\noindent\textbf{4. (C4 - Menganalisis \textbar\ Bobot: 20 Poin)}\par\vspace{0.5mm}
{\footnotesize Berdasarkan fungsi arus $i(t)$ pada Soal 3, analisis profil disipasi daya sesaat pada resistor $p_R(t) = R i^2(t)$ dan laju penyimpanan energi medan magnetik pada induktor $w_L(t) = \frac{1}{2} L i^2(t)$. Hitung energi total yang tersimpan dalam induktor saat sistem telah mencapai keadaan tunak (\textit{steady-state}, $t \to \infty$).}
\begin{answerbox}{Ruang Pengerjaan C4 (Analisis Neraca Energi \& Daya Transien)}
  \vspace{3.8cm}
\end{answerbox}

\vspace{1.5mm}

% Soal C5
\noindent\textbf{5. (C5 - Mengevaluasi \textbar\ Bobot: 20 Poin)}\par\vspace{0.5mm}
{\footnotesize Suatu koil relay kontaktor gardu memiliki spesifikasi $R = 12\,\Omega$ dan $L = 2{,}4\text{ H}$. Pabrikan menyatakan bahwa arus kontaktor membutuhkan waktu minimal $3\tau$ untuk mencapai batas arus angkat (*pick-up current* $95\%$ dari $I_{\text{maks}}$). Evaluasi pernyataan teknis tersebut secara matematis, dan buktikan apakah waktu $t = 3\tau$ memang mencukupi batas $95\%$ tersebut!}
\begin{answerbox}{Ruang Pengerjaan C5 (Evaluasi Kinerja Teknis Kontaktor)}
  \vspace{3.8cm}
\end{answerbox}

\vspace{1.5mm}

% Soal C6
\noindent\textbf{6. (C6 - Merancang \& Komputasi Julia \textbar\ Bobot: 15 Poin)}\par\vspace{0.5mm}
{\footnotesize Rancang skrip program ilmiah berbasis \textbf{Julia} (\texttt{Differential\-Equations.jl} \& \texttt{Plots.jl}) untuk menyelesaikan PDB $\frac{di}{dt} = \frac{V_s - R i}{L}$ dengan solver \texttt{Tsit5()} pada selang $t \in [0, 0{,}5]\text{ s}$. Tuliskan sintaks fungsi ODE dan blok pemecahannya di bawah ini.}
\begin{answerbox}{Ruang Pengerjaan C6 (Sintaks Program Ilmiah Julia \& Parameter Desain)}
  \vspace{3.8cm}
\end{answerbox}

\vspace{1.5mm}

% Tabel Rekapitulasi Skor OBE & Tanda Tangan
\noindent\begin{minipage}[b]{0.65\textwidth}%
  \centering
  \scriptsize
  \begin{tabular}{|c|c|c|c|c|c|c|}
    \hline
    \rowcolor{unibblue}
    \textbf{\color{white}C1} & \textbf{\color{white}C2} & \textbf{\color{white}C3} & \textbf{\color{white}C4} & \textbf{\color{white}C5} & \textbf{\color{white}C6} & \textbf{\color{white}Total} \\
    \rowcolor{unibblue!10}
    10 Pts & 15 Pts & 20 Pts & 20 Pts & 20 Pts & 15 Pts & 100 Pts \\
    \hline
    & & & & & & \\[2.5mm]
    \hline
  \end{tabular}%
\end{minipage}%
\hfill%
\begin{minipage}[b]{0.32\textwidth}%
  \centering
  \scriptsize
  \textbf{Verifikasi Dosen / Asisten:}\\[5mm]
  (\dotfill)
\end{minipage}\par

\end{document}
